\documentclass[10pt,a4paper]{amsart}
\usepackage[margin=19mm]{geometry}
\usepackage{amsmath,amssymb,mathtools}
\usepackage[scr=rsfs,cal=euler]{mathalfa}
\usepackage{xfrac}
\usepackage[unicode,hidelinks,bookmarks=false]{hyperref}
\newcommand{\ie}{i.e.\ }
\newcommand{\cf}{cf.\ }
\newcommand{\resp}{respectively\ }
\newcommand{\Subsection}{\subsection}
\providecommand{\nobreakdash}{\nobreak\hbox{-}}
\setlength{\emergencystretch}{2em}
\multlinegap0pt
\usepackage{amssymb}
\usepackage[shortlabels]{enumitem}
\newtheorem{theorem}{Theorem}
\newtheorem{proposition}{Proposition}
\newtheorem{lemma}{Lemma}
\DeclareMathOperator{\Imm}{Imm}
\DeclareMathOperator{\Ort}{O}
\DeclareMathOperator{\SO}{SO}
\DeclareMathOperator{\diam}{diam}
\DeclareMathOperator{\dist}{dist}
\DeclareMathOperator{\e}{e}
\DeclareMathOperator{\lspan}{span}
\DeclareMathOperator*{\osc}{osc}
\DeclareMathOperator{\vol}{vol}
\title{\textbf{Asymptotic rigidity of codimension-1 isometric immersions via quantitative estimates}}
\author{\textbf{Mert Baştuğ}}
\date{Source: arXiv:2604.09387v1 (CC BY 4.0)}
\begin{document}
\maketitle

\noindent\emph{Source and licence.} \textbf{Asymptotic rigidity of codimension-1 isometric immersions via quantitative estimates}. Version arXiv:2604.09387v1 (CC BY 4.0), distributed by arXiv under the Creative Commons Attribution 4.0 International licence, http://creativecommons.org/licenses/by/4.0/, as recorded in the arXiv licence metadata for this exact version and shown on the abstract page https://arxiv.org/abs/2604.09387v1. Full licence notice: \texttt{licenses/arxiv-2604.09387-CC-BY-4.0.txt}.\par\smallskip

\noindent\textit{Editorial scope.} Counted unit: the article's theorem thm:codimension1\_compact\_rigidity (source0.tex lines 161--167). The article's complete original proof of it is delivered with it (source0.tex lines 406--461, one \texttt{proof} environment of 56 lines, which the article places away from the statement). No mathematics is written, repaired or re-derived: the statement and the proof are the article's own lines, in the article's own notation, and no retained line is substituted or re-flowed.  Retained uncounted as the source-local context the proof needs: lines 226--232; lines 258--268; lines 295--300; lines 313--315; lines 317--319.  Nothing was omitted from any retained span, so the package carries no omission record.  No retained line contains a citation, so the wrapper carries no bibliography and no bibliography entry.  No retained line contains a figure environment, an image-inclusion command or drawing code, so no image is delivered and the package declares no asset dependency.  Editorial formatting only, in addition: an AMS \texttt{amsart} preamble that restates the article's own declarations and macros used by the retained lines, namely the environments \texttt{theorem}, \texttt{proposition}, \texttt{lemma} on separate counters, together with the standard packages those lines need.  The retained environments print in this document as Theorem 1, Proposition 1, Lemma 1 in order of appearance, whereas the article prints the same items with the numbering of its own full document.  The complete native source of the article as distributed in the arXiv e-print for this exact version is bundled unchanged as \texttt{sources/198204/source0.tex}.
\begin{theorem} \label{thm:noneuclidean_rigidity}
    For every $u \in W^{1, p}(Q; \mathbb{R}^d)$ and $x_0 \in Q$, there exists $R \in \SO(g_{x_0}, \e_d)$ such that
    \[
    \int_Q |Du - R|_{g, \e_d}^p \, dx \le C\left(|Q|\left(\osc_Q g\right)^p + \int_Q \dist_{g, \e_d}^p(Du, \SO(g, \e_d)) \, dx\right).
    \]
    The constant $C$ depends only on $p$, $d$ and $\lambda$.
\end{theorem}

\begin{proposition} \label{pr:derivative_normal_estimate}
    Let $u \in \Imm_p(Q; N)$. Then for all $j = 1, \dots, d$,
    \begin{equation} \label{eq:derivative_of_the_normal_bound1}
        |\partial_{x_j} \bar{\nu}_u|^2 \le |P_u(\partial_{x_j} \bar{\nu}_u)|^2 + C |\partial_{x_j} \bar{u}|^2 \quad \text{a.e. in } Q.
    \end{equation}
    In particular,
    \[
    \int_Q |D\bar{\nu}_u|_{g, \e_D}^p \, dx \le C \mathcal{E}(u).
    \]
    The constant $C$ depends only on $d$, $p$, $\lambda$, $N$ and $\iota$.
\end{proposition}

\begin{lemma} \label{lm:change_in_volume}
    If $f : Q \rightarrow [0, \infty]$ is measurable, then
    \[
    \frac{1}{\lambda^{d/2}} \int_Q f \, dx \le \int_Q f \, d\vol_g \le \lambda^{d/2} \int_Q f \, dx.
    \]
\end{lemma}

    \begin{equation} \label{eq:error_in_projection}
        |P_{\Pi_0}T - T|_{g_x, \e_D} \le |T|_{g_x, \e_D}|\Pi_0^\perp - \Pi^\perp| \quad \text{for all } x \in Q.
    \end{equation}

    \begin{equation} \label{eq:bound_for_distance_from_oriented_rotations}
        \dist_{g_x, e_D}\left(P_{\Pi_0} T, \SO\left((\mathbb{R}^d, g_x), (\Pi_0, \e_D)\right)\right) \le \dist_{g_x, e_D}(T, \Ort(g_x, \e_D)) + C|\Pi_0^\perp - \Pi^\perp| \quad \text{for all } x \in Q,
    \end{equation}

\begin{theorem} \label{thm:codimension1_compact_rigidity}
    Let $u \in \Imm_p(Q; N)$. There exist $x_0 \in Q$ and $R \in \Ort(g_{x_0}, \e_D)$ such that
    \[
    \int_Q |D\bar{u} - R|_{g, \e_D}^p \, dx \le C \left(|Q| \left(\osc_Q g\right)^p + E_s(u) + \diam^p(Q) \mathcal{E}(u)\right).
    \]
    The constant $C$ depends only on $d$, $p$, $\lambda$, $N$ and $\iota$.
\end{theorem}

\begin{proof}[Proof of Theorem \ref{thm:codimension1_compact_rigidity}]
    Roughly speaking, our strategy will be to project $\bar{u}$ to a linear space and to apply the non-Euclidean rigidity theorem \ref{thm:noneuclidean_rigidity}. In order to estimate the error due to the projection, we need to bound the variation of the tangent spaces $D\bar{u}(x)(\mathbb{R}^d)$ as $x$ ranges over $Q$.
    
    We shall denote $D\bar{u}(x)(\mathbb{R}^d)$ simply by $\Pi_x$. The subspace of $\mathbb{R}^D$ canonically isomorphic to $T_q\iota(N)$ is denoted by $V_q$. We claim that
    \begin{equation} \label{eq:normal_plane_variation}
        \int_Q \int_Q |\Pi_x^\perp - \Pi_y^\perp|^p \, dx \, dy \le C \diam^p(Q) |Q| \mathcal{E}(u).
    \end{equation}
    Since $\Pi_x^\perp = \lspan(\bar{\nu}_u(x)) \oplus V_{\bar{u}(x)}^\perp$ for a.e. $x \in Q$, we have
    \[
    |\Pi_x^\perp - \Pi_y^\perp|^2 \le |\bar{\nu}_u(x) - \bar{\nu}_u(y)|^2 + |V_{\bar{u}(x)}^\perp - V_{\bar{u}(y)}^\perp|^2 \quad \text{for a.e. } x, y \in Q.
    \]
    Set $r := D - d - 1$. Let $U \subset \mathbb{R}^D$ be a convex open set and let $n_i \in C^\infty(U; \mathbb{R}^D)$ for $i = 1, \dots, r$ such that $(n_1(q), \dots, n_r(q))$ is a positively oriented orthonormal basis of $V_{q}^\perp$ for all $q \in U \cap \iota(N)$, and $|Dn_i|$ is uniformly bounded in $U$ for all $i$. For $q, q' \in U \cap \iota(N)$, we have
    \[
    |V_q^\perp - V_{q'}^\perp|^2 \le \sum_{i = 1}^r |n_i(q) - n_i(q')|^2 \le Cr |q - q'|^2,
    \]
    where $C$ is a uniform bound for $|Dn_i|$ in $U$. Since $\iota(N)$ is compact, a covering argument yields
    \[
    |V_q^\perp - V_{q'}^\perp| \le C |q - q'| \quad \text{for all } q, q' \in \iota(N)
    \]
    with $C$ depending on $d$, $D$, $N$ and $\iota$. By Poincaré's inequality,
    \begin{equation} \label{eq:planar_variation}
        \int_Q \int_Q |V_{\bar{u}(x)}^\perp - V_{\bar{u}(y)}^\perp|^p \, dx \, dy \le C \int_Q \int_Q |\bar{u}(x) - \bar{u}(y)|^p \, dx \, dy \le C \diam^p(Q) |Q| \int_Q |D\bar{u}|^p \, dx.
    \end{equation}
    Applying Poincaré's inequality to $\bar{\nu}_u$ and using Proposition \ref{pr:derivative_normal_estimate} gives
    \begin{equation} \label{eq:normal_variation}
        \int_Q \int_Q |\bar{\nu}_u(x) - \bar{\nu}_u(y)|^p \, dx \, dy \le C \diam^p(Q) |Q| \int_Q |D\bar{\nu}_u|^p \, dx \le C \diam^p(Q) |Q| \mathcal{E}(u).
    \end{equation}
    The bound \eqref{eq:normal_plane_variation} now follows from \eqref{eq:planar_variation} and \eqref{eq:normal_variation} with the help of Lemma \ref{lm:change_in_volume}.

    By Chebyshev's inequality applied to \eqref{eq:normal_plane_variation}, there exists $x_0 \in Q$ with $\dim \Pi_{x_0} = d$ such that
    \begin{equation} \label{eq:poincare_oriented_distance}
        \int_Q |\Pi_x^\perp - \Pi_{x_0}^\perp|^p \, dx \le 2C \diam^p(Q) \mathcal{E}(u).
    \end{equation}
    Let $T \in \SO((\mathbb{R}^d, \e_d), (\Pi_{x_0}, \e_D))$ and set $v := T^{-1} \circ P_{\Pi_{x_0}} \circ \bar{u}$. Then $v \in W^{1, p}(Q; \mathbb{R}^d)$, and we can apply Theorem \ref{thm:noneuclidean_rigidity} to obtain $\bar{R} \in \SO(g_{x_0}, \e_d)$ such that
    \begin{equation} \label{eq:rigidity_applied_to_v_again}
        \int_Q |Dv - \bar{R}|_{g, \e_d}^p \, dx \le C \left(|Q| \left(\osc_Q g\right)^p + \int_Q \dist_{g, \e_d}^p(Dv, \SO(g, \e_d)) \, dx\right).
    \end{equation}
    Set $R := T\bar{R}$. Then $R \in \SO((\mathbb{R}^d, g_{x_0}), (\Pi_{x_0}, \e_D))$, and by \eqref{eq:error_in_projection}, we have
    \begin{equation} \label{eq:change_from_v_to_u}
        \begin{aligned}
            |D\bar{u}(x) - R|_{g_x, \e_D} &\le |D\bar{u}(x) - P_{\Pi_{x_0}} D\bar{u}(x)|_{g_x, \e_D} + |P_{\Pi_{x_0}} D\bar{u}(x) - R|_{g_x, \e_D} \\
            &\le |\Pi_x^\perp - \Pi_{x_0}^\perp| |D\bar{u}(x)|_{g_x, \e_D} + |Dv(x) - \bar{R}|_{g_x, \e_d} \\
            &\le \sqrt{d} |\Pi_x^\perp - \Pi_{x_0}^\perp| + 2\sqrt{D - d} \dist_{g_x, \e_D}(D\bar{u}(x), \Ort(g, \e_D)) + |Dv(x) - \bar{R}|_{g_x, \e_d}
        \end{aligned}
    \end{equation}
    for a.e. $x \in Q$. Furthermore, \eqref{eq:bound_for_distance_from_oriented_rotations} gives
    \begin{multline} \label{eq:planar_to_spacial_distance}
        \dist_{g_x, \e_d}(Dv(x), \SO(g_x, \e_d)) = \dist_{g_x, \e_D}(P_{\Pi_{x_0}} D\bar{u}(x), \SO((\mathbb{R}^d, g_x), (\Pi_{x_0}, \e_D)) \\
        \le \dist_{g_x, \e_D}(D\bar{u}(x), \Ort(g_x, \e_D)) + C|\Pi_x^\perp - \Pi_{x_0}^\perp|
    \end{multline}
    for a.e. $x \in Q$. Hence, \eqref{eq:rigidity_applied_to_v_again}, \eqref{eq:change_from_v_to_u}, \eqref{eq:planar_to_spacial_distance} and Lemma \ref{lm:change_in_volume} imply
    \[
    \int_Q |D\bar{u} - R|_{g, \e_D}^p \, dx \le C \left(|Q| \left(\osc_Q g\right)^p + E_s(u) + \int_Q |\Pi_x^\perp - \Pi_{x_0}^\perp|^p \, dx\right).
    \]
    Finally, the claim follows from \eqref{eq:poincare_oriented_distance}.
\end{proof}

\end{document}
